\begin{table}[!htbp]
\centering\small
\caption{跨候选关联及共同按天重采样区间。}\label{tab:corr}
\begin{tabularx}{\linewidth}{@{}lXrcrc@{}}
\toprule
组 & 横轴指标 & $r$ & 95\% 区间 & $\rho$ & 95\% 区间 \\
\midrule
v1 & RMSE & 0.351 & [$-$0.171, 0.535] & 0.482 & [$-$0.303, 0.638] \\
v1 & 日中心 RMSE & 0.252 & [$-$0.152, 0.382] & 0.391 & [$-$0.238, 0.559] \\
v1 & Listwise 代理 & 0.320 & [$-$0.241, 0.485] & 0.374 & [$-$0.338, 0.553] \\
v1 & 窗口排序 $T$ & 0.849 & [0.472, 0.881] & 0.612 & [0.200, 0.812] \\
v2 & RMSE & 0.005 & [$-$0.438, 0.333] & 0.082 & [$-$0.474, 0.321] \\
v2 & 日中心 RMSE & $-$0.108 & [$-$0.488, 0.203] & $-$0.035 & [$-$0.484, 0.233] \\
v2 & Listwise 代理 & $-$0.165 & [$-$0.518, 0.124] & $-$0.039 & [$-$0.460, 0.093] \\
v2 & 窗口排序 $T$ & 0.777 & [0.197, 0.852] & 0.656 & [0.135, 0.814] \\
\bottomrule
\end{tabularx}
\par\vspace{3pt}{\footnotesize $r$ 为 Pearson；$\rho$ 为跨候选 Spearman。窗口 $T$ 本身另含窗口内 Spearman。7 日块、10,000 次 percentile 重采样，仅改变日期、固定候选集合；区间不包含候选搜索或抽样不确定性。}
\end{table}
